\section{Very Large Circles: Base 2}

To use the formula $J(n)=2\ell+1$, we must first write $n=2^m+\ell$, that is,
find the largest power of $2$ not bigger than $n$. For a large $n$ this is not
immediate in base $10$. In base $2$ it is: the decomposition can be read
directly on the digits of $n$, and the formula becomes a simple operation on
these digits. In this whole section, the game starts with person $1$, as in
Section~\ref{sec:formula}: $J(n)$ is the survivor of the usual game.

\subsection{Reading \texorpdfstring{$2^m$}{2\^{}m} and \texorpdfstring{$\ell$}{l} in base 2}\label{sec:bindecomp}

In base $2$, a number is written with the digits $0$ and $1$, and each digit
stands for a power of $2$. For instance
\[
10=(1010)_2=1\cdot2^3+0\cdot2^2+1\cdot2^1+0\cdot2^0=8+2 .
\]
Let $n\geq1$, with $2^m\leq n<2^{m+1}$. Then $n$ has exactly $m+1$ digits in
base $2$, and the first one is a $1$:
\[
n=(1\,b_{m-1}\,b_{m-2}\cdots b_0)_2
 =2^m+\underbrace{b_{m-1}2^{m-1}+\dots+b_12+b_0}_{=\,\ell}.
\]
Since each $b_i$ is $0$ or $1$, the remainder satisfies
$0\leq\ell\leq2^{m-1}+\dots+2+1=2^m-1<2^m$. So this is exactly the
decomposition of Observation~\ref{obs:decomp}: the first digit gives $2^m$,
and the other digits give $\ell$. For $n=10=(1\,010)_2$ we read $2^3=8$ and
$\ell=(010)_2=2$.

\subsection{The rule: move the first digit to the end}

\begin{theo}\label{thm:binary}
If the game starts with person $1$ and $n=(1\,b_{m-1}\cdots b_1b_0)_2$, then
\[
J(n)=(b_{m-1}\cdots b_1b_0\,1)_2 .
\]
In words: to get $J(n)$, we take the first digit of $n$ in base $2$, which is
a $1$, and move it to the end.
\end{theo}

\begin{proof}
By Section~\ref{sec:bindecomp}, $\ell=(b_{m-1}\cdots b_0)_2=\sum_{i=0}^{m-1}b_i2^i$, and
$J(n)=2\ell+1$ by Theorem~\ref{thm:main}. Multiplying by $2$ moves every digit
one place to the left, and adding $1$ writes a $1$ in the last place:
\[
2\ell+1=\sum_{i=0}^{m-1}b_i2^{i+1}+1
=b_{m-1}2^m+b_{m-2}2^{m-1}+\dots+b_02^1+1\cdot2^0
=(b_{m-1}\cdots b_0\,1)_2 .
\]
\end{proof}

Table~\ref{tab:binary} shows the rule on our sixteen tests. When the digit
after the first one is $0$, the result starts with $0$, and this $0$ is simply
dropped: for $n=8=(1000)_2$ we get $(0001)_2=1$.

\begin{table}[!ht]
\centering
\renewcommand{\arraystretch}{1.1}
\begin{tabular}{|c|c|c|c|}
\hline
$n$ & $n$ in base 2 & $J(n)$ in base 2 & $J(n)$\\ \hline
1 & 1 & 1 & 1\\
2 & 10 & 01 & 1\\
3 & 11 & 11 & 3\\
4 & 100 & 001 & 1\\
5 & 101 & 011 & 3\\
6 & 110 & 101 & 5\\
7 & 111 & 111 & 7\\
8 & 1000 & 0001 & 1\\
9 & 1001 & 0011 & 3\\
10 & 1010 & 0101 & 5\\
11 & 1011 & 0111 & 7\\
12 & 1100 & 1001 & 9\\
13 & 1101 & 1011 & 11\\
14 & 1110 & 1101 & 13\\
15 & 1111 & 1111 & 15\\
16 & 10000 & 00001 & 1\\ \hline
\end{tabular}
\caption{The first digit of $n$ moves to the end. The values agree with
Table~\ref{tab:results}.}
\label{tab:binary}
\end{table}

\subsection{What base 2 shows}

Base $2$ gives another view of Corollary~\ref{cor:consequences}.
\begin{itemize}
  \item \textbf{$J(n)$ and $n$ have the same number of digits $1$.} Moving a
        digit, or dropping a $0$ at the front, does not change the number of
        $1$'s. For example $10=(1010)_2$ and $J(10)=5=(101)_2$ both have two.
  \item \textbf{$J(n)=1$ exactly when $n$ is a power of $2$.} The result
        $(b_{m-1}\cdots b_0\,1)_2$ is equal to $1$ only when all the $b_i$ are
        $0$, that is when $n=(10\cdots0)_2=2^m$.
  \item \textbf{$J(n)=n$ exactly when all the digits of $n$ are $1$.} The
        strings $1\,b_{m-1}\cdots b_1b_0$ and $b_{m-1}\cdots b_1b_0\,1$ are equal
        when $b_{m-1}=1$, $b_{m-2}=b_{m-1}$, \dots, $b_0=b_1$, so when all the
        digits are $1$: $n=(11\cdots1)_2=2^{m+1}-1$, as for $n=3$, $7$, $15$.
\end{itemize}

\Needspace{14\baselineskip}
\subsection{Computing \texorpdfstring{$J(n)$}{J(n)} by hand for a large \texorpdfstring{$n$}{n}}

To find $J(n)$ by hand for a large $n$, we follow four steps:
\begin{enumerate}
  \item write $n$ in base $2$, by dividing by $2$ again and again and keeping
        the remainders;
  \item remove the first digit, which is a $1$;
  \item write a $1$ at the end of what is left;
  \item convert the result back to base $10$.
\end{enumerate}

\textbf{Example: $n=1000$.} We divide by $2$ until we reach $0$:
\begin{center}
\renewcommand{\arraystretch}{1.1}
\begin{tabular}{l*{10}{c}}
number & 1000 & 500 & 250 & 125 & 62 & 31 & 15 & 7 & 3 & 1\\
remainder & 0 & 0 & 0 & 1 & 0 & 1 & 1 & 1 & 1 & 1
\end{tabular}
\end{center}
Reading the remainders from right to left gives $1000=(1111101000)_2$. We move
the first $1$ to the end:
\[
(1\,111101000)_2\ \longrightarrow\ (111101000\,1)_2=512+256+128+64+16+1=977 .
\]
So $J(1000)=977$. With the formula: $1000=512+488$, so
$J(1000)=2\times488+1=977$.

\textbf{Example: $n=2026$.} We have $2026=(11111101010)_2$. Moving the first digit to the end gives
\[
(1\,1111101010)_2\ \longrightarrow\ (1111101010\,1)_2=2005,
\]
so $J(2026)=2005$. This agrees with
$2026=2^{10}+1002$ and $2\times1002+1=2005$.

\textbf{Larger values.} Table~\ref{tab:large} gives a few larger values, all obtained by moving the
first digit. We also checked them with the program of Section~6. In the line
$n=10\,000$, two zeros appear at the front after the move and are dropped.

\noindent\begin{minipage}{\textwidth}\centering
\renewcommand{\arraystretch}{1.15}
\footnotesize\setlength{\tabcolsep}{4pt}
\begin{tabular}{|r|l|r|l|}
\hline
$n$ & $n$ in base 2 & $J(n)$ & $J(n)$ in base 2\\ \hline
$1000$ & 1111101000 & $977$ & 1111010001\\
$2026$ & 11111101010 & $2005$ & 11111010101\\
$10\,000$ & 10011100010000 & $3617$ & 111000100001\\
$10^6$ & 11110100001001000000 & $951\,425$ & 11101000010010000001\\
$2^{20}$ & 1 followed by twenty 0 & $1$ & 1\\
$2^{20}-1$ & twenty digits 1 & $1\,048\,575$ & twenty digits 1\\
$2^{20}+1$ & 1, nineteen 0, 1 & $3$ & 11\\
$10^9$ & 111011100110101100101000000000 & $926\,258\,177$ & 110111001101011001010000000001\\ \hline
\end{tabular}
\captionof{table}{Large values of $J(n)$. In each line, $n$ and $J(n)$ have the same
number of digits $1$.}
\label{tab:large}
\end{minipage}
\par\medskip

For $n=10^6$, there are $20$ digits, so $m=19$, $2^{19}=524\,288$ and
$\ell=10^6-524\,288=475\,712$, which gives $J(10^6)=951\,425$. For $n=10^9$,
the game itself needs $999\,999\,999$ eliminations. With base
$2$, we only handle the $30$ digits of $n$. In general, $n$ has about
$\log_2 n$ digits in base $2$, so the work grows like $\log_2 n$ instead of
$n$.
